\documentclass[UTF8,12pt]{ctexart}
\usepackage[a4paper,margin=24mm]{geometry}
\usepackage{amsmath,amssymb,bm,graphicx,adjustbox}
\newcommand{\fitmath}[1]{\begin{adjustbox}{max width=\linewidth}$\displaystyle #1$\end{adjustbox}}
\setlength{\parindent}{0pt}
\setlength{\parskip}{.7em}
\sloppy
\allowdisplaybreaks
\begin{document}
\section*{2019 年考研数学三 · 参考答案}
\subsection*{原 PDF 第 5 页}

\subsection*{一、选择题}

（1）C。　（2）D。　（3）D。　（4）B。　（5）A。　（6）C。　（7）C。　（8）A。


\subsection*{二、填空题}

（9）\(\displaystyle e^{-1}\)。　（10）\(\displaystyle (\pi,\allowbreak -2)\)。　（11）\(\displaystyle \dfrac{\displaystyle 1-2\sqrt2}{\displaystyle 18}\)。　（12）0.4。　（13）1。　（14）\(\displaystyle \dfrac{\displaystyle 2}{\displaystyle 3}\)。


\subsection*{三、解答题}

（15）\(\displaystyle f^\prime(x)=\begin{cases}e^x(x+1),\allowbreak &x<0,\allowbreak \\2e^{2x\ln x}(\ln x+1),\allowbreak &x>0.\end{cases}\)


\(\displaystyle x=-1\) 和 \(\displaystyle x=\dfrac{\displaystyle 1}{\displaystyle e}\) 是 \(\displaystyle f(x)\) 的极小值点，极小值分别为 \(\displaystyle f(-1)=1-e^{-1}\) 和 \(\displaystyle f(\dfrac{\displaystyle 1}{\displaystyle e})=e^{-2/e}\)；\(\displaystyle x=0\) 是 \(\displaystyle f(x)\) 的极大值点，极大值为 \(\displaystyle f(0)=1\)。


（16）\(\displaystyle 1-3f_{11}^{\prime\prime}(x+y,\allowbreak x-y)-f_{22}^{\prime\prime}(x+y,\allowbreak x-y)\)。


（17）（I）\(\displaystyle y(x)=\sqrt x\,e^{x^2/2}\)。（II）\(\displaystyle \dfrac{\displaystyle 1}{\displaystyle 2}\pi(e^4-e)\)。


（18）\(\displaystyle \dfrac{\displaystyle 1}{\displaystyle 2}+\dfrac{\displaystyle 1}{\displaystyle e^\pi-1}\)。


（19）（I）证明略。（II）\(\displaystyle \lim_{n\to\infty}\dfrac{\displaystyle a_n}{\displaystyle a_{n-1}}=1\)。


（20）当 \(\displaystyle a\neq-1\) 时，向量组 I 与向量组 II 等价。当 \(\displaystyle a=1\) 时，\(\displaystyle \boldsymbol\beta_3=(3-2k)\boldsymbol\alpha_1+(-2+k)\boldsymbol\alpha_2+k\boldsymbol\alpha_3\)，其中 \(\displaystyle k\) 为任意常数；当 \(\displaystyle a\neq\pm1\) 时，\(\displaystyle \boldsymbol\beta_3=\boldsymbol\alpha_1-\boldsymbol\alpha_2+\boldsymbol\alpha_3\)。


（21）（I）\(\displaystyle x=3,\allowbreak y=-2\)。（II）满足 \(\displaystyle \boldsymbol P^{-1}\boldsymbol A\boldsymbol P=\boldsymbol B\) 的可逆矩阵为


\[\fitmath{\displaystyle \boldsymbol P=\begin{pmatrix}-1&-1&-1\\2&1&2\\0&0&4\end{pmatrix}.}\]


（22）（I）\(\displaystyle Z\) 的概率密度为 \(\displaystyle f_Z(z)=\begin{cases}pe^z,\allowbreak &z\leq0,\allowbreak \\(1-p)e^{-z},\allowbreak &z>0.\end{cases}\)


（II）当 \(\displaystyle p=\dfrac{\displaystyle 1}{\displaystyle 2}\) 时，\(\displaystyle X\) 与 \(\displaystyle Z\) 不相关。（III）\(\displaystyle X\) 与 \(\displaystyle Z\) 不相互独立。


（23）（I）\(\displaystyle A=\sqrt{\dfrac{\displaystyle 2}{\displaystyle \pi}}\)。（II）\(\displaystyle \sigma^2\) 的最大似然估计量为 \(\displaystyle \widehat{\sigma}^2=\dfrac{\displaystyle \sum_{i=1}^n(X_i-\mu)^2}{\displaystyle n}\)。


\clearpage
\end{document}
